\textbf{Purpose:} Patients are often asked to prove a simple fact about
their health record, such as having received at least a required number of
vaccine doses, and end up disclosing the whole record. A zero-knowledge
proof can answer the question without revealing the data, but only if the
proof is tied to a stored record and to a verification session. We address
this for a patient-held record stored on IPFS and anchored on Solana.

\textbf{Methods:} Starting from a prototype in which a server computed the
proof from the private value, we store a salted Poseidon commitment to the
value in the record's on-chain account, make the same commitment an output
of the circuit, add a verifier-chosen challenge as a public input, enforce
input ranges inside the circuit, and generate Groth16 proofs in the owner's
browser. We implement the design with Circom, snarkjs and Anchor, test it,
and measure it on a laptop.

\textbf{Results:} The bound circuit has \NewConstraints{} constraints
against \OldConstraints{} for the baseline. A proof takes a median of
\BrowserProveMs~ms to generate in a browser and \NewVerifyMs~ms to verify.
Each record occupies a \AccountBytes-byte account with a refundable deposit
of \RentSol{} SOL, and its creation consumes at most \CuMax{} compute units.

\textbf{Conclusion:} Binding a threshold proof to an on-chain commitment and
proving on the owner's device is practical for this claim. The committed
value remains self-asserted; an issuer attestation is the main missing
piece.
